\documentclass[11pt]{article}

% ---- portable preamble: only standard CTAN packages, compiles offline ----
\usepackage[utf8]{inputenc}
\usepackage[T1]{fontenc}
\usepackage[a4paper,margin=1in]{geometry}
\usepackage{amsmath,amssymb}
\usepackage{booktabs}
\usepackage{array}
\usepackage{multirow}
\usepackage{pifont}
\usepackage{graphicx}
\usepackage{xcolor}
\usepackage[hidelinks]{hyperref}
\usepackage[round]{natbib}
\bibliographystyle{plainnat}

\newcommand{\sys}{\textsc{synthitd}}
\newcommand{\code}[1]{\texttt{\small #1}}

\title{\textbf{\sys: A Reproducible, Endogenous-Intent Synthetic Benchmark\\
for Insider-Threat Detection}}

\author{
  Ahmed Alaoui Mdaghri\thanks{Correspondence: \href{mailto:mad.mik788@gmail.com}{mad.mik788@gmail.com}}\\
  \texttt{LLM\_DataGen} project
}
\date{\today}

\begin{document}
\maketitle

\begin{abstract}
Insider-threat detection (ITD) research is bottlenecked by data: real incidents are
rare and sensitive, so a decade of work trains and evaluates almost entirely on the
synthetic CERT corpus. Systematic review shows the resulting scores are inflated by
balanced subsets, random splits, and a hand-scripted malicious population whose
\emph{signature} a model can fingerprint instead of learning intent. Recent
LLM-based successors improve realism (ChimeraLog) and verifiability (OrgForge-IT) but
still \emph{inject} a scripted attack into a chosen user and ship neither a benign
anomaly population nor a self-audit. We present \sys, a reproducible generator and an
honest benchmark that targets these shortcomings. Its central idea is to
\emph{grow} insider intent from a latent Critical-Pathway risk model
(predisposition $\rightarrow$ stressors $\rightarrow$ concerning behaviour) with a
hazard that decides whether and when a pathway activates, rather than injecting a
script---so labels track latent intent, not a template. A deterministic, label-owning
event bus produces multi-surface telemetry (logon, IdP, device, file, email, chat,
ticketing, version control, HTTP, DLP, endpoint, HR) while a \emph{label-blind}
language-model renderer produces free text that cannot leak the answer. The dataset
ships graded labels (benign / anomalous-benign / precursor / malicious), an engineered
benign-anomaly hard-negative population, a sensor-degradation model, and a CERT-r6.2
compatible exporter. \sys{} also ships its own benchmark with temporal,
user-held-out, and \emph{scenario-held-out} splits, operational metrics
(AUC-PR at true prevalence, detection earliness, review-budget recall, false-positive
attribution), and a leakage/shortcut audit. We show the single after-hours$+$removable
$+$upload rule that nears ceiling on CERT collapses to PR-AUC $0.0004$ on \sys{}, and
that a faithful reimplementation of the best-reported CERT architecture
(UBS-Transformer; reported AUROC $0.95$/F1 $0.96$) drops to AUROC $0.64$/F1 $0.43$ on
\sys. Everything is a deterministic function of a seed and regenerates for free.
\end{abstract}

\section{Introduction}
Malicious insiders cause some of the costliest security incidents, yet detection
research is constrained less by modelling than by \emph{data}. Real traces are scarce,
privacy-laden, and legally fraught, so the field has standardised on the synthetic
CERT Insider Threat Test Dataset~\citep{glasser2013bridging}. A decade of results built
on CERT report near-perfect numbers that do not transfer to operations
~\citep{jurisic2026decade}: the malicious users were generated from a handful of
templates, so detectors learn the \emph{generator's} fingerprint---an after-hours
logon followed by a removable-media copy and an upload---rather than intent. Reported
performance is further inflated by evaluating on class-balanced subsets and random
(non-temporal) splits.

Two 2026 LLM-based datasets advance the state of the art. ChimeraLog
~\citep{yu2026chimera} uses multi-agent LLMs to raise realism to near that of a real
competition dataset, and shows classifiers that score $\sim$0.99 on CERT fall to
$\sim$0.83 on it. OrgForge-IT~\citep{orgforgeit2026} makes ground truth an
architectural guarantee by letting a deterministic engine own a \code{SimEvent} bus
while LLMs render only prose. Both, however, still \emph{inject} a scripted attack
into a chosen user and label at the granularity of the injected steps; neither ships a
constructed benign-anomaly population or an audit of how much of its own signal is a
generator artifact. MOLE~\citep{mole2026} reframes the problem for AI agents and
confirms the field's move toward budget-aware, review-cost metrics.

\paragraph{Contributions.}
\begin{enumerate}
  \item \textbf{Endogenous intent.} Each employee carries a latent Critical-Pathway
  risk trajectory~\citep{shaw2015critical,sofit2018}; a hazard model decides whether
  and when a pathway activates. Many predisposed, stressed employees never activate and
  remain hard negatives. The label tracks latent intent, defeating signature shortcuts.
  \item \textbf{Deterministic bus $+$ label-blind renderer.} A typed event bus owns
  identities, timestamps, causal links, and the graded label; an LLM renders only free
  text and never sees the label, enforced by a leakage linter.
  \item \textbf{Engineered hard negatives, graded labels, sensor model.} A benign-but-
  anomalous population (travel, on-call, promotion, approved bulk export) produces
  exfiltration-looking bursts with justifying context; labels are graded
  (benign / anomalous-benign / precursor / malicious); an observability layer models
  partial sensor coverage.
  \item \textbf{An honest benchmark} with temporal, user-held-out, and
  scenario-held-out splits; AUC-PR at true prevalence; detection earliness;
  review-budget recall; false-positive attribution; and a shipped shortcut audit.
  \item \textbf{Comparison to the best-reported CERT SOTA.} We faithfully reimplement
  the UBS-Transformer~\citep{ubstransformer2025} and show its reported CERT performance
  does not transfer to \sys. A CERT-r6.2 compatible exporter makes \sys{} a drop-in for
  existing pipelines.
\end{enumerate}

This is a defensive, academic artifact: it models behaviour as detection-relevant
observables and MITRE ATT\&CK technique \emph{labels}~\citep{ctid2022insider}, not
operational tradecraft.

\section{Background and Gap Analysis}
\label{sec:gap}
Table~\ref{tab:gap} situates \sys{} against the incumbent and its recent successors.
CERT~\citep{glasser2013bridging} composes a topic model, a social graph, per-user
OCEAN psychometrics and an LDAP role model to emit logon/device/file/email/HTTP logs;
its authors are explicit that realism is bounded and malicious scenarios are
hand-scripted and few. The canonical releases are r4.2 (1{,}000 users, dense), r5.2
(2{,}000, 99 insiders), and r6.2 (4{,}000, 5 insiders, $\sim$0.1\% positives). The
honest supervised SOTA on the full imbalanced population is a random forest on
user-day features at ROC-AUC $\approx 0.979$ \citep{le2020granularity}; the inflated
end reports $\geq 0.99$ on balanced/SMOTE subsets \citep{sivakrishna2025xgboost}.

\begin{table}[t]
\centering
\caption{Capability comparison. \sys{} adds endogenous intent, graded labels, an
engineered hard-negative population, a sensor model, and a shipped shortcut audit.}
\label{tab:gap}
\small
\begin{tabular}{lcccc}
\toprule
Capability & CERT r6.2 & ChimeraLog & OrgForge-IT & \textbf{\sys} \\
\midrule
Deterministic label-owning bus      & \ding{55} & $\sim$ & \checkmark & \checkmark \\
Modern surfaces (chat/ticket/VCS/IdP)& \ding{55} & partial & \checkmark & \checkmark \\
\textbf{Endogenous intent}          & \ding{55} & \ding{55} (injected) & \ding{55} (injected) & \textbf{\checkmark} \\
\textbf{Graded labels}              & \ding{55} & \ding{55} & partial & \textbf{\checkmark} \\
\textbf{Benign-anomaly hard negatives} & \ding{55} & \ding{55} & partial & \textbf{\checkmark} \\
\textbf{Sensor-degradation model}   & \ding{55} & \ding{55} & \ding{55} & \textbf{\checkmark} \\
Regeneratable at any size/prevalence& \ding{55} & expensive & \checkmark & \checkmark \\
\textbf{Leakage / shortcut audit}   & \ding{55} & \ding{55} & \ding{55} & \textbf{\checkmark} \\
Earliness \& review-budget metrics  & \ding{55} & \ding{55} & partial & \textbf{\checkmark} \\
CERT-compatible export              & n/a & \ding{55} & \ding{55} & \textbf{\checkmark} \\
\bottomrule
\end{tabular}
\end{table}

\section{Generator Design}
\label{sec:design}
\sys{} is a deterministic function of a configuration (organisation size, horizon,
prevalence, noise, seed). A single seed feeds a named-substream RNG, so independent
concerns draw independently and the dataset is stable under unrelated code changes.

\paragraph{Organisation.} Three data-sensitive domains (technology, finance,
healthcare) define roles with privilege levels and sensitive-asset access. Each
employee receives Big-Five (OCEAN) traits and a static \emph{predisposition}. Insiders
are sampled super-linearly in predisposition (floored so anyone may become one), and a
separate benign-anomaly population is assigned.

\paragraph{Critical-Pathway psychology (endogenous intent).}
For each employee and day $d$ we integrate a latent risk trajectory
\begin{equation}
r_d = \sigma\!\big(-1.9 + 2.4\,\pi + 1.3\,s_d + 1.1\,c_d + 1.6\,\iota_d + 0.6\,\delta_d\big),
\end{equation}
where $\pi$ is the static predisposition, $s_d$ an exponentially-decaying accumulation
of life/HR stressors (negative review, passed-over promotion, sanction, financial
hardship, resignation), $c_d$ a concerning-behaviour feedback term, $\iota_d$ a
\emph{latent} intent-pressure ramp present only for the insider population, and
$\delta_d$ autocorrelated drift. $\iota_d$ is never emitted as an observable feature,
so it creates no shortcut; normal employees can still reach high risk through
stressors, so high risk $\neq$ insider. A hazard converts elevated risk into a per-day
activation probability; the first success is the onset day, and many insiders never
activate within the horizon.

\paragraph{Behaviour and hard negatives.} A baseline model emits per-day activity
across twelve surfaces, modulated by role and personality. Legitimate power users
(and, more heavily, the engineered benign-anomaly population) upload to the cloud, use
removable media, and work after hours, so those features are \emph{not} markers of
malice. Benign triggers (travel, on-call, promotion, crunch, approved bulk export)
produce exfiltration-looking bursts \emph{paired with an observable HR context event}
that explains them, making precision and false-positive attribution measurable.

\paragraph{Threat pathways.} Five abstract, stage-based playbooks (IP theft, data
leak, privilege misuse, fraud, sabotage markers) fire only for activated insiders.
Each stage emits the same typed events as the baseline, differing in target, volume,
timing, and label, and carries MITRE ATT\&CK technique-id labels (15 distinct
techniques). A per-insider \emph{stealth} level (from conscientiousness and the
risk-curve shape) spreads covert actors into work hours, small volumes, and
benign-looking channels, while variant sampling randomises target, channel, and
staging order---so the positive class is a distribution, not a template.

\paragraph{Observability and label-blind rendering.} An observability layer applies
per-channel coverage and dropout to produce the \emph{observed} stream a detector
sees, retaining the full stream for evaluation and reporting how much malicious signal
was rendered invisible. Free-text surfaces are produced by a separate pass given only a
label-free \code{ContentPlan}; a leakage linter redacts any label term or technique id.
Consequently the renderer writes ordinary prose for malicious events too, forcing
detectors onto behavioural signal rather than keyword matching.

\section{Dataset}
\label{sec:dataset}
Every observable is an \code{Event} with a channel, actor, action, structured
attributes, a graded \code{Label}, optional MITRE technique ids, the latent risk at
event time, a causal parent, and an optional rendered-text payload. Labels are graded:
\emph{benign}; \emph{anomalous\_benign} (the hard-negative class, paired with a
justifying HR event); \emph{precursor} (concerning activity on an armed pathway, before
the act); and \emph{malicious}. Table~\ref{tab:stats} reports a shipped sample and a
scaled reference instance; both regenerate deterministically from their config.

\begin{table}[t]
\centering
\caption{Two shipped instances. Prevalence is per malicious \emph{user-day}; the
reference is sparser than CERT r6.2.}
\label{tab:stats}
\small
\begin{tabular}{lrr}
\toprule
 & Sample (tech) & Reference (finance) \\
\midrule
Employees                      & 120      & 1{,}000 \\
Horizon (days)                 & 70       & 120 \\
Full events                    & 797{,}328 & 9{,}132{,}134 \\
Observed events                & 705{,}641 & 7{,}827{,}513 \\
Insider population / activated  & 7 / 5    & 8 / 6 \\
Anomalous-benign employees     & 10       & 80 \\
Benign events                  & 792{,}555 & 9{,}107{,}771 \\
Anomalous-benign events        & 4{,}560  & 24{,}103 \\
Precursor events               & 139      & 199 \\
Malicious events               & 74       & 61 \\
Malicious user-day prevalence  & --       & 0.02\% \\
\bottomrule
\end{tabular}
\end{table}

\section{Benchmark}
\label{sec:benchmark}
The unit of analysis is the (user, day). Features are built only from the observed
stream and only from attributes a real sensor could report; generator-internal flags
that encode the label are excluded by an explicit denylist, and a regression test
asserts that injecting them changes no feature. The observable HR-context event
\emph{is} a feature, so a good model can learn to exonerate explained anomalies.

\paragraph{Splits.} \emph{Temporal} (train early days, test later) prevents the
future-into-past leakage of random splits; \emph{user-held-out} tests generalisation
to unseen people; \emph{scenario-held-out} removes a whole pathway family from
training---a generalisation test prior datasets do not offer.

\paragraph{Metrics.} We report ROC-AUC and AUC-PR at the true sparse prevalence (not on
a balanced subset), precision/recall at a review budget (top $b\%$ of user-days),
detection \emph{earliness} (days between first flag and onset; enabled by precursor
labels), and false-positive attribution to the anomalous-benign class.

\paragraph{Shortcut audit.} We ship an audit that scores a single trivial rule
(after-hours $+$ removable $+$ upload) against a full model. On CERT that rule is
near-perfect; on the \sys{} reference it reaches PR-AUC $0.0004$ (ROC $0.547$) while a
full logistic model reaches PR-AUC $0.092$---a positive \emph{signal-depth gap} of
$+0.091$, and the rule's false positives are $64\%$ the engineered hard negatives.
The shortcut is removed by construction.

\section{Experiments}
\label{sec:exp}

\subsection{Reference baselines}
Table~\ref{tab:baselines} runs six reference detectors on the finance reference
($120{,}000$ user-days, $0.02\%$ malicious). At this prevalence the trivial rule is
useless and static detectors are weak in PR-AUC; the causal temporal model
(\code{sequence\_logreg}) is the only one that stays strong on ranking, generalises to
unseen users and the held-out \code{data\_leak} scenario, and---under a $1\%$ review
budget---catches far more insiders, with a median detection lead of up to two weeks.

\begin{table}[t]
\centering
\caption{Reference baselines (finance, $1000\times120$, $0.02\%$ prevalence).
R@1\%: recall at a $1\%$ review budget. Lead: median days before the act (positive =
early). FP$_{ab}$: share of flagged negatives that are engineered hard negatives.}
\label{tab:baselines}
\small
\setlength{\tabcolsep}{4.5pt}
\begin{tabular}{llrrrrr}
\toprule
Split & Detector & ROC & PR-AUC & R@1\% & Lead & FP$_{ab}$ \\
\midrule
\multirow{3}{*}{temporal}
 & volume\_rule (CERT shortcut) & 0.547 & 0.0004 & 0.00 & -- & 0.64 \\
 & logistic regression          & 0.809 & 0.092 & 0.42 & $-1$ & 0.15 \\
 & \textbf{sequence\_logreg}    & \textbf{0.991} & \textbf{0.478} & \textbf{0.83} & $0$ & 0.17 \\
\midrule
\multirow{3}{*}{user-holdout}
 & volume\_rule                 & 0.738 & 0.0008 & 0.00 & 36 & 0.56 \\
 & logistic regression          & 0.829 & 0.099 & 0.57 & 30 & 0.14 \\
 & \textbf{sequence\_logreg}    & \textbf{1.000} & \textbf{0.613} & \textbf{1.00} & 16.5 & 0.22 \\
\midrule
\multirow{3}{*}{scenario-holdout}
 & volume\_rule                 & 0.553 & 0.0003 & 0.00 & -- & 0.70 \\
 & logistic regression          & 0.773 & 0.103 & 0.38 & 14 & 0.08 \\
 & \textbf{sequence\_logreg}    & \textbf{0.996} & \textbf{0.441} & \textbf{0.88} & 14 & 0.10 \\
\bottomrule
\end{tabular}
\end{table}

\subsection{SOTA recipe family, ablated}
To attribute the difficulty to \sys{}'s mechanisms rather than to a confound, we
generate two datasets that share seed, org, onset, labels, prevalence and splits and
differ only in a \code{legacy\_cert\_mode} flag that degenerately disables the whole
anti-shortcut regime (benign power-user uploads/removable, benign after-hours baseline,
benign-anomaly bursts, and insider stealth). We run reference reimplementations of the
dominant CERT recipes over five seeds (Table~\ref{tab:ablation}). Every classical
detector drops $43$--$92\%$ in PR-AUC and $\sim$0.1 ROC from the CERT regime to the
realistic regime; on realistic data all are weak in PR-AUC against a no-skill baseline
of $\approx0.0017$. We report ROC and PR together and do not cross metrics: the
classical collapse holds on both, while the temporal model's advantage is on ranking
and budgeted recall, not PR.

\begin{table}[t]
\centering
\caption{SOTA recipe families, CERT-regime $\rightarrow$ realistic (temporal split,
mean over 5 seeds). The whole anti-shortcut regime is toggled at once; see
\S\ref{sec:threats}.}
\label{tab:ablation}
\small
\begin{tabular}{lcccc}
\toprule
Detector (recipe family) & ROC (CERT$\to$real) & PR-AUC (CERT$\to$real) & $\Delta$PR \\
\midrule
Random forest (user-day)        & $0.879\to0.773$ & $0.686\to0.363$ & $-47\%$ \\
Gradient boosting (XGBoost-like)& $0.860\to0.610$ & $0.598\to0.049$ & $-92\%$ \\
MLP (deep feed-forward)         & $0.825\to0.652$ & $0.642\to0.341$ & $-47\%$ \\
Autoencoder (PCA reconstruction)& $0.880\to0.817$ & $0.679\to0.281$ & $-59\%$ \\
Temporal GBDT (sequence)        & $0.966\to0.879$ & $0.445\to0.111$ & $-75\%$ \\
\bottomrule
\end{tabular}
\end{table}

\subsection{The best-reported CERT SOTA, reimplemented}
We take the strongest recent CERT architecture with a fully specified design---the
UBS-Transformer~\citep{ubstransformer2025}: user-based sequencing, a
6-layer/$d_{\text{model}}{=}512$/8-head reconstruction Transformer encoder trained with
MSE on normal users, then One-Class SVM / LOF / Isolation Forest on the per-user
reconstruction errors, evaluated per user. It reports accuracy $0.966$, precision
$0.935$, recall $0.994$, F1 $0.964$, AUROC $0.950$ on CERT. We reimplement it exactly
(adapting the per-step token to \sys{}'s day-granularity user-day features) and run it
on a CERT-regime version of \sys{} (a faithfulness check) and on realistic \sys{}
(Table~\ref{tab:ubs}). On CERT-regime data the reimplementation reproduces the paper's
regime (AUROC $0.87$--$0.89$), landing slightly lower as expected for a stylized
in-pipeline proxy trained on far less data and fewer epochs. The \emph{same code at the
same scale} then drops to AUROC $0.643$ / F1 $0.425$ on realistic data; the bare
reconstruction signal falls to AUROC $0.573$, barely above chance. The best-reported
CERT SOTA loses $\sim$0.23 AUROC and $\sim$0.26 F1 and does not transfer.

\begin{table}[t]
\centering
\caption{UBS-Transformer~\citep{ubstransformer2025}: reported on CERT vs. our faithful
reimplementation on \sys{} ($400\times120$, per-user evaluation). The
within-reimplementation CERT-regime $\rightarrow$ realistic gap is the confound-free
measurement; the paper row is the external anchor.}
\label{tab:ubs}
\small
\begin{tabular}{llrrrr}
\toprule
Setting & Outlier det. & AUROC & F1 & Precision & Recall \\
\midrule
Paper, reported on CERT        & iForest & \textbf{0.950} & \textbf{0.964} & 0.935 & 0.994 \\
Reimpl, CERT-regime \sys{}     & iForest & 0.874 & 0.681 & 0.571 & 0.842 \\
Reimpl, CERT-regime \sys{}     & OCSVM   & 0.892 & 0.750 & 0.714 & 0.789 \\
Reimpl, realistic \sys{}       & iForest & \textbf{0.643} & \textbf{0.425} & 0.357 & 0.526 \\
Reimpl, realistic \sys{}       & LOF     & 0.703 & 0.516 & 0.667 & 0.421 \\
\bottomrule
\end{tabular}
\end{table}

\section{Threats to Validity}
\label{sec:threats}
We state limitations plainly. (i) \emph{Model-based, not learned from real logs:}
baseline rates are hand-set and realism is asserted via stylized-facts checks, not
validated against a proprietary corpus. (ii) \emph{The \code{legacy\_cert\_mode}
ablation toggles four coupled axes at once} (including insider stealth, which
attenuates the positive signal), so Table~\ref{tab:ablation} reflects the combined
regime; the isolated signature effect is the separate shortcut audit. (iii)
\emph{Under-power:} experiments run at tens of insiders and $\sim$0.17\% prevalence, so
reported magnitudes are indicative and the direction is the result; the SOTA ablation
is reported over five seeds, the UBS reimplementation over one. (iv) \emph{Metric
honesty:} at these prevalences AUC-PR is near a $\sim$0.002 no-skill floor, so all
realistic detectors are operationally weak; we do not claim any model ``works''. (v)
\emph{Proxies:} the PCA autoencoder is a lower bound on the DNN/LSTM auto-encoder
family, the temporal GBDT is a proxy for LSTM/transformer sequence models, and the
CERT regime is reproduced in-pipeline rather than measured on the real corpus (which
was not available in our environment).

\section{Ethics, Availability, and Reproducibility}
\sys{} is fully synthetic: no real people, accounts, credentials, or messages; names
are drawn from a fixed neutral list. It models behaviour as detection observables and
MITRE technique labels and contains no operational attack instructions; it is intended
for defensive and academic use only. Every dataset is a pure function of its
configuration and regenerates byte-for-byte; the deterministic template renderer runs
at zero API cost, with a Claude-based renderer as an optional, label-blind quality
upgrade. The code, configs, benchmark, audits, and this paper are released together.

\section{Conclusion}
\sys{} replaces injected scripts with grown intent, guarantees ground truth through a
label-owning bus and a label-blind renderer, ships engineered hard negatives, graded
labels, a sensor model, and a self-critical benchmark. The signature shortcut that
inflates CERT is removed by construction, and the best-reported CERT SOTA does not
transfer to it---evidence that realistic insider-threat detection remains unsolved and
that the useful signal lives in the behavioural trajectory, not a static signature.

\small
\bibliography{references}

\end{document}
